% STATE: DRAFT (awaiting human review).
\section{Background and Design}\label{sec:design}

\subsection{Preliminaries}

% PROV[user]: Current request; prior discussion of labelled Kripke executions.
% PROV[code]: KILT/LTL/Syntax.lean:L3-L30; KILT/LTL/Semantics.lean:L4-L65;
% KILT/Kripke/Path.lean; KILT/Kripke/Semantics.lean.
KILT uses the following core grammar for future-time LTL over atomic
propositions $p \in AP$:
\[
  \varphi ::= \bot \mid p \mid \neg\varphi
    \mid \varphi \lor \varphi
    \mid \mathsf{X}\varphi
    \mid \varphi\,\mathsf{U}\,\varphi .
\]
Conjunction, implication, eventually, globally and release are derived
operators.  A Kripke execution $s_0s_1\ldots$ induces the infinite valuation
word $L(s_0)L(s_1)\ldots$, over which satisfaction is defined at a position.
This follows the standard presentation in Baier and
Katoen~\cite{BaierKatoen2008}; the separation of syntax from trace semantics
was also informed by LeanLTL~\cite{VinMillerFremont2025}.

% PROV[user]: Current request; prior explanations of Büchi acceptance and
% finite-lasso counterexamples.
% PROV[code]: KILT/Graph/Acceptance.lean:L24-L35,L81-L86;
% KILT/Graph/Lasso.lean:L1-L12.
For a B\"uchi automaton $A=(Q,\Sigma,\delta,Q_0,F)$, a run
$\rho=q_0q_1\ldots$ is accepting when
\[
  \operatorname{Inf}(\rho) \cap F \neq \varnothing,
\]
where $\operatorname{Inf}(\rho)$ is the set of states occurring infinitely
often.  Generalised B\"uchi acceptance replaces $F$ by
$\{F_1,\ldots,F_n\}$ and requires
$\operatorname{Inf}(\rho)\cap F_i\neq\varnothing$ for every $i$.  In a finite
graph this reduces to a reachable accepting cycle, represented by a finite
prefix and a non-empty loop---a \emph{lasso}---repeated forever.

\subsection{Overall Architecture}

% PROV[user]: Current request and prior pipeline summary.
% PROV[code]: KILT/LTL/Compile.lean:L3-L34,L41-L82;
% KILT/Automata/Product.lean:L21-L29,L49-L98; KILT/Kripke/Check.lean:L17-L46;
% KILT/Kripke/Counterexample.lean:L25-L64.
To check a model against a property $\varphi$, KILT compiles $\lnot\varphi$.
The formula is normalised into negation normal form and expanded into a tableau
whose states are consistent assignments to its subformulas.  One recurring
acceptance condition for each strong-until obligation turns this tableau into a
generalised B\"uchi automaton; a verified round-robin counter then produces an
ordinary finite B\"uchi automaton.  Its synchronous product with the Kripke
structure tracks both the current system state and the automaton's progress on
the system label.  A graph search for a reachable accepting cycle finally
returns either no violation or a counterexample lasso.  The final checker
theorem connects absence of such a cycle with satisfaction by every initial
infinite execution.

\subsection{Design Choices and Reused Libraries}

% PROV[user]: Current request; prior observation that correctness is layered
% and that Graph, automata and tableau components should remain reusable.
% PROV[code]: KILT/LTL/Compile.lean:L6-L34;
% KILT/Automata/GeneralizedBuchiCorrect.lean; KILT/Graph/Acceptance.lean:L5-L16;
% lake-manifest.json (pinned CSLib dependency).
I kept the construction modular: declarative LTL semantics, formula
compilation, automata-language preservation, product construction and finite
graph search have separate interfaces and correspondence theorems.  This makes
``correctness'' a chain of precise claims rather than one monolithic assertion.
CSLib supplies omega sequences and the underlying nondeterministic and B\"uchi
run semantics~\cite{CSLib}; KILT adds the executable finite automata, explicit
tableau, generalised-to-ordinary B\"uchi conversion, product, lasso search and
their connecting proofs.  The tableau follows the obligation-based idea of
Gerth, Peled, Vardi and Wolper~\cite{GerthPeledVardiWolper1995}, but is built
explicitly rather than on demand.  Keeping the graph layer independent of LTL
also isolates the finite reasoning that makes both safety and liveness checking
executable.
